\chapter{Background and Related Work}
\label{chap:background}

\section{Transaction fraud as imbalanced learning}

Fraud detection differs from a balanced classroom classification task. Fraud is
rare, labels can be delayed, behavior changes, and the cost of a missed fraud
case may differ substantially from the cost of reviewing a legitimate
transaction. Accuracy can therefore be misleading: a classifier that predicts
the majority class can appear accurate while detecting no fraud.

For a positive fraud class, precision and recall are:

\begin{align}
\mathrm{Precision} &= \frac{TP}{TP + FP},\\
\mathrm{Recall} &= \frac{TP}{TP + FN}.
\end{align}

ROC-AUC measures ranking across true- and false-positive rates, while PR-AUC
summarizes the precision--recall tradeoff. Davis and Goadrich showed the
relationship between ROC and PR spaces and highlighted the usefulness of PR
curves under class skew \parencite{davis2006pr}. Saito and Rehmsmeier further
demonstrated why PR plots can be more informative for imbalanced binary
classifiers \parencite{saito2015pr}.

\section{The IEEE-CIS dataset}

The source consists of transaction tables and smaller optional identity tables.
The public columns include transaction amount, product and card categories,
email domains, device attributes, time, and many anonymized groups such as
\texttt{C}, \texttt{D}, \texttt{M}, and \texttt{V}
\parencite{ieeecis2019,amazonfraudbenchmark}. The training transaction file
contains 590,540 rows in the inspected source, and the natural fraud prevalence
is approximately 3.5\%.

The project uses a left join because identity information is unavailable for
most transactions. Missingness and identity presence can themselves be useful
signals, but they must be interpreted carefully: association does not show that
missing data causes fraud.

\section{Apache Spark for batch processing}

Apache Spark provides a unified engine and higher-level APIs for distributed
data processing \parencite{zaharia2016spark}. This project uses Spark for typed
CSV ingestion, joins, profiling, feature computation, chronological splitting,
Parquet export, and a small decision-tree demonstration.

Spark is not used for request-time inference. Model training crosses from Spark
to pandas because the compared Python estimators use in-memory array/data-frame
interfaces. This boundary simplifies integration but constrains training scale
to available driver memory.

\section{Gradient-boosted decision trees}

Gradient boosting builds an additive ensemble of decision trees, where later
trees correct residual errors of the current ensemble
\parencite{friedman2001greedy}. Three modern implementations are compared:

\begin{description}[style=multiline,leftmargin=2.6cm,labelwidth=2.3cm]
  \item[XGBoost] provides scalable tree boosting with sparsity-aware learning,
  weighted quantile sketches, cache-aware access, and distributed execution
  \parencite{chen2016xgboost}.
  \item[LightGBM] uses histogram-based learning together with techniques such as
  gradient-based one-side sampling and exclusive feature bundling
  \parencite{ke2017lightgbm}.
  \item[CatBoost] introduces ordered boosting and categorical-statistic
  techniques designed to reduce prediction shift and target leakage
  \parencite{prokhorenkova2018catboost}.
\end{description}

Logistic regression provides a lower-complexity baseline. The comparison uses
the same training/validation interface, but the final V1/V2 comparison is not a
controlled algorithm-only experiment because feature contracts differ.

\section{Explainability with SHAP}

SHAP is a framework for additive feature attribution based on Shapley values
\parencite{lundberg2017shap}. For a prediction \(f(x)\), an additive explanation
approximates the output as:

\begin{equation}
g(z') = \phi_0 + \sum_{i=1}^{M}\phi_i z'_i,
\end{equation}

where \(\phi_i\) is the contribution attributed to feature \(i\). The project
uses a tree-specific explainer and returns the five largest absolute local
contributions.

SHAP values provide a model explanation, not a causal explanation. They can
also explain a different output scale from the probability shown to a user when
a separate calibrator transforms the raw model score.

\section{Service and interface technologies}

FastAPI uses Python type annotations and Pydantic models to create runtime
validation and OpenAPI schemas \parencite{fastapi}. SQLAlchemy maps the
transaction and review entities to PostgreSQL or SQLite
\parencite{sqlalchemy,postgresql}. React provides the component model for the
browser application \parencite{react}, while Vite builds and serves the
TypeScript bundle \parencite{vite}.

The technologies are connected as a modular monorepo rather than independent
microservices. This is appropriate for a short project because it reduces
deployment complexity, but the backend image inherits heavyweight training
dependencies from the model package.

\section{Related engineering concerns}

A production fraud platform would additionally require:

\begin{itemize}
  \item a feature store or request-time method that exactly reproduces
  historical aggregates;
  \item monitored calibration and drift;
  \item business-approved operating thresholds;
  \item immutable model lineage and controlled promotion;
  \item durable case management and review audit history;
  \item authentication, authorization, secrets, and privacy governance.
\end{itemize}

The present work evaluates these as readiness gaps instead of claiming that the
academic demonstration already provides them.
